
\subsubsection{Key Findings}

The central finding is that computational feasibility validation was absent from this automated pipeline's workflow between design and implementation phases. Specific observations:

\begin{enumerate}
\item \textbf{Orthogonality of feasibility and quality:} 100\% implementation task completion, passing tests, and SDD compliance coexisted with complete execution infeasibility. Traditional quality metrics do not detect resource constraints.
\end{enumerate}

\begin{enumerate}
\item \textbf{Late discovery cost:} Infeasibility was discovered after 10-16 hours of implementation effort (planning + coding). Early detection at Phase 2C.5 (~5 minutes) could have prevented this waste in this case.
\end{enumerate}

\begin{enumerate}
\item \textbf{Estimation complexity:} Naive calculation (model + optimizer = 282GB) substantially underestimated actual requirements (489GB). The gap (207GB, 73\% underestimate) suggests intuition-based feasibility judgments may be unreliable for large models.
\end{enumerate}

\begin{enumerate}
\item \textbf{Pipeline gap:} The YouRA automated pipeline lacked a checkpoint between experiment design (Phase 2C) and implementation (Phase 3-4) to validate resource feasibility. Phase 2C focused on scientific requirements without hardware constraint validation; Phase 3-4 assumed approved designs were executable.
\end{enumerate}

\subsubsection{Limitations}

\textbf{Single case study:} This analysis documents one failure in one automated pipeline. We cannot claim the feasibility gate prevents failures in diverse workflows without broader testing. A single case suffices to demonstrate that such gaps can exist and to propose a solution design, but validating gate effectiveness across multiple projects is necessary future work.

\textbf{Gate not yet validated:} The proposed gate is a design based on failure analysis, not a validated tool with empirical effectiveness data. We provide retrospective validation on one case (correct prediction) but cannot guarantee it prevents future failures without false positives (rejecting feasible configurations) or false negatives (missing infeasible ones). Prospective deployment across diverse projects is required.

\textbf{Estimation formula approximations:} The formula uses conservative approximations for activation memory and framework overhead (~10\%). Actual values vary with runtime factors (dynamic batch sizes, gradient checkpointing, framework versions). Conservative estimates may produce false positives. The 85\% utilization threshold provides safety margin but requires empirical calibration.

\textbf{Cost-benefit assumes compliance:} The 120:1 to 192:1 ratio assumes researchers would heed warnings and reformulate, with negligible false positive costs. In practice, researchers might override warnings, false positives could block feasible experiments, and overly conservative gates might deter ambitious experiments. Actual value depends on gate accuracy, researcher behavior, and organizational culture.

\textbf{Narrow scope:} Findings come from one automated pipeline (YouRA) with specific characteristics (automated workflow, MoE Transformer, multi-task NLP, memory constraints only). Manual workflows may have informal feasibility practices not captured here. Different domains may have different constraint profiles. The paper title and scope reflect this: we identify a gap in one automated system, not universal workflows.

\subsubsection{Implications}

\textbf{For automated research pipelines:} This case suggests automated systems could benefit from systematic feasibility validation between design and implementation phases. While experienced practitioners may perform informal checks, automated pipelines lack these validations. As models scale beyond single-device capacity, such gaps may result in more frequent failures.

\textbf{For resource-constrained research:} Labs with fixed hardware budgets may find value in avoiding wasted limited resources. Early validation enables informed resource allocation decisions.

\textbf{For reproducibility:} Infeasibility failures often go unreported. Making constraints explicit and documented improves transparency about what experiments are runnable versus aspirational.

\subsubsection{Future Work}

Immediate extensions:
\begin{enumerate}
\item \textbf{Multi-project validation:} Deploy Phase 2C.5 gate across 20-30 research projects spanning different model scales, tasks, and hardware to measure false positive/negative rates empirically
\item \textbf{Framework-specific calibration:} Collect overhead measurements for PyTorch, JAX, TensorFlow to refine estimation accuracy
\item \textbf{Threshold tuning:} Investigate whether 85\% utilization should vary based on multi-GPU count, framework, or experimental complexity
\end{enumerate}

Longer-term directions:
\begin{enumerate}
\item \textbf{Automated feasibility tools:} Integrate estimation formulas into experiment design interfaces
\item \textbf{Resource-aware design:} Tools that suggest feasible configurations given hardware constraints
\item \textbf{Community benchmarking:} Public database of memory requirements for popular models across frameworks
\item \textbf{Multi-constraint checking:} Extend beyond memory to dataset availability, time budgets, license compliance
\end{enumerate}
